\section{Rozszerzenia Kana}\label{cat:sec:kan-extensions}
Rozszerzenia Kana uogólniają konstrukcję granic, kogranic i sprzężeń: zamiast rozszerzać funkcję z podzbioru, poszukujemy uniwersalnego przedłużenia funktora wzdłuż innego funktora, który nie musi być włożeniem. Wykorzystamy kategorie przecinkowe z definicji~\ref{cat:def:comma}, granice i kogranice z rozdziału~\ref{sec:granice}, transformacje naturalne z rozdziału~\ref{sec:naturalne} oraz sprzężenia z rozdziału~\ref{sec:adjunctions}. W całym rozdziale $\mathcal A,\mathcal B$ są małymi kategoriami, $\mathcal C$ jest lokalnie mała. Kategorie funktorów oznaczamy przez $[\mathcal A,\mathcal C]$ i $[\mathcal B,\mathcal C]$; ich zbiory transformacji naturalnych są podzbiorami małych iloczynów zbiorów morfizmów.

\subsection{Własności uniwersalne}
\begin{definicja}[Lewe i prawe rozszerzenie Kana]\label{cat:def:kan-extension}
Niech $K:\mathcal A\to\mathcal B$ i $F:\mathcal A\to\mathcal C$. Lewym rozszerzeniem Kana funktora $F$ wzdłuż $K$ nazywamy parę $(L,\eta)$ z $L:\mathcal B\to\mathcal C$ i transformacją $\eta:F\Rightarrow LK$, dla której dla każdego $H:\mathcal B\to\mathcal C$ odwzorowanie
\[
 \Nat(L,H)\longrightarrow\Nat(F,HK),\qquad
 \alpha\longmapsto(\alpha K)\eta
\]
jest bijekcją. Piszemy $L=\operatorname{Lan}_K F$.
Prawym rozszerzeniem Kana nazywamy parę $(R,\epsilon)$ z $R:\mathcal B\to\mathcal C$ i $\epsilon:RK\Rightarrow F$, dla której dla każdego $H$ odwzorowanie
\[
 \Nat(H,R)\longrightarrow\Nat(HK,F),\qquad
 \beta\longmapsto\epsilon(\beta K)
\]
jest bijekcją. Piszemy $R=\operatorname{Ran}_K F$. Iloczyn transformacji oznacza ich składanie pionowe; $(\alpha K)_a=\alpha_{K(a)}$.
\end{definicja}
\begin{center}
\begin{tabular}{cc}
\(\begin{array}{ccc}
\mathcal A&\xrightarrow{K}&\mathcal B\\
&\searrow F&\downarrow L\\
&&\mathcal C
\end{array}\)
&
\(\begin{array}{ccc}
\mathcal A&\xrightarrow{K}&\mathcal B\\
&\searrow F&\downarrow R\\
&&\mathcal C
\end{array}\)\\
\(\eta:F\Rightarrow LK\)&\(\epsilon:RK\Rightarrow F\)
\end{tabular}
\end{center}
Nie wymagamy $LK=F$ ani nawet izomorficzności tych funktorów. Warunek uniwersalny dotyczy transformacji, nie równości funktorów.

\begin{stwierdzenie}[Jednoznaczność]\label{cat:prop:kan-unique}
Dwa lewe rozszerzenia Kana $(L,\eta),(L',\eta')$ są połączone jedynym izomorfizmem naturalnym $\theta:L\Rightarrow L'$ spełniającym $(\theta K)\eta=\eta'$. Dla prawych rozszerzeń analogicznie istnieje jedyny izomorfizm $\theta:R\Rightarrow R'$ z $\epsilon'(\theta K)=\epsilon$.
\end{stwierdzenie}
\begin{proof}
Z własności uniwersalnej $L$ otrzymujemy jedyną $\theta$ z podanym równaniem, a z własności $L'$ jedyną $\gamma:L'\Rightarrow L$ z $(\gamma K)\eta'=\eta$. Wtedy $((\gamma\theta)K)\eta=(\gamma K)\eta'=\eta$, więc jednoznaczność daje $\gamma\theta=\id_L$. Analogicznie $((\theta\gamma)K)\eta'=\eta'$ daje $\theta\gamma=\id_{L'}$. W przypadku prawym własność $R'$ daje $\theta$, a własność $R$ daje $\gamma:R'\Rightarrow R$ z $\epsilon(\gamma K)=\epsilon'$. Mamy $\epsilon((\gamma\theta)K)=\epsilon'(\theta K)=\epsilon$ i $\epsilon'((\theta\gamma)K)=\epsilon(\gamma K)=\epsilon'$, więc oba złożenia są identycznościami.
\end{proof}
\begin{twierdzenie}[Rozszerzenia jako sprzężenia]\label{cat:thm:kan-adjunction}
Niech $K^*:[\mathcal B,\mathcal C]\to[\mathcal A,\mathcal C]$ oznacza prekompozycję z $K$. Jeżeli lewe rozszerzenia istnieją dla wszystkich $F$, wyznaczają funktor $\operatorname{Lan}_K$ lewy sprzężony do $K^*$. Jeżeli prawe rozszerzenia istnieją dla wszystkich $F$, wyznaczają funktor $\operatorname{Ran}_K$ prawy sprzężony do $K^*$. Odwrotnie, istnienie odpowiedniego sprzężenia daje rozszerzenia Kana.
\end{twierdzenie}
\begin{proof}
Wybieramy rozszerzenia $(L_F,\eta_F)$. Dla $\mu:F\Rightarrow F'$ definiujemy $L_\mu:L_F\Rightarrow L_{F'}$ jako jedyną transformację z
\[
 (L_\mu K)\eta_F=\eta_{F'}\mu.
\]
Dla identyczności równanie spełnia $\id_{L_F}$. Dla $\nu:F'\Rightarrow F''$ mamy $((L_\nu L_\mu)K)\eta_F=(L_\nu K)\eta_{F'}\mu=\eta_{F''}\nu\mu$, więc jednoznaczność daje zachowanie składania. Własność uniwersalna daje bijekcję $\Nat(L_F,H)\cong\Nat(F,K^*H)$. Naturalność względem $H$ wynika z łączności składania, a względem $F$ z równania definiującego $L_\mu$: obraz $\alpha L_\mu$ jest równy $(\alpha K)\eta_{F'}\mu$. Jest to sprzężenie w sensie definicji~\ref{def:adjunction}.

Dla prawych rozszerzeń $(R_F,\epsilon_F)$ definiujemy $R_\mu:R_F\Rightarrow R_{F'}$ przez $\epsilon_{F'}(R_\mu K)=\mu\epsilon_F$. Identyczności spełniają to równanie, a dla $\nu$ mamy $\epsilon_{F''}((R_\nu R_\mu)K)=\nu\mu\epsilon_F$, co daje zachowanie składania. Bijekcja $\Nat(H,R_F)\cong\Nat(K^*H,F)$ jest naturalna względem $H$ przez prekompozycję, a względem $F$ przez równanie $\epsilon_{F'}(R_\mu K)=\mu\epsilon_F$. Odwrotnie jednostka sprzężenia $\operatorname{Lan}_K\dashv K^*$ daje $\eta_F$, a kojednostka $K^*\dashv\operatorname{Ran}_K$ daje $\epsilon_F$; bijekcje sprzężenia są dokładnie bijekcjami z definicji~\ref{cat:def:kan-extension}.
\end{proof}

\subsection{Wzory punktowe}
\begin{definicja}\label{cat:def:kan-comma}
Dla $b\in\Ob(\mathcal B)$ używamy dwóch kategorii przecinkowych. W $(K\downarrow b)$ obiektem jest $(a,u:K(a)\to b)$, a morfizmem do $(a',u')$ jest $f:a\to a'$ z $u'K(f)=u$. W $(b\downarrow K)$ obiektem jest $(a,v:b\to K(a))$, a morfizmem do $(a',v')$ jest $f:a\to a'$ z $K(f)v=v'$. Projekcje $P_b^-,P_b^+$ na $\mathcal A$ zapominają strzałkę $u$ lub $v$. Obie kategorie są małe: ich obiekty tworzą sumę rozłączną zbiorów morfizmów indeksowaną obiektami $\mathcal A$, a morfizmy są ograniczone warunkiem przemienności.
\end{definicja}
\begin{twierdzenie}[Punktowe lewe rozszerzenie Kana]\label{cat:thm:pointwise-lan}
Jeżeli dla każdego $b$ istnieje kogranica diagramu $FP_b^-:(K\downarrow b)\to\mathcal C$, to istnieje lewe rozszerzenie Kana i można je obliczyć przez
\[
 (\operatorname{Lan}_K F)(b)=\colim_{(a,u:K(a)\to b)} F(a).
\]
W szczególności warunek zachodzi, gdy $\mathcal C$ jest kozupełna.
\end{twierdzenie}
\begin{proof}
Wybieramy kogranicę $L(b)$ i morfizmy jej kostożka $\lambda^b_{a,u}:F(a)\to L(b)$, zatem $\lambda^b_{a',u'}F(f)=\lambda^b_{a,u}$ dla morfizmów kategorii przecinkowej. Dla $g:b\to c$ rodzina $\lambda^c_{a,gu}$ spełnia to samo równanie, gdyż z $u'K(f)=u$ wynika $(gu')K(f)=gu$. Własność uniwersalna kogranicy wyznacza jedyny $L(g)$ z
\begin{equation}\label{cat:eq:lan-action}
 L(g)\lambda^b_{a,u}=\lambda^c_{a,gu}.
\end{equation}
Dla $g=\id_b$ równanie spełnia $\id_{L(b)}$; dla $h:c\to d$ obie strony $L(hg)$ i $L(h)L(g)$ dają po prekompozycji z każdą $\lambda^b_{a,u}$ morfizm $\lambda^d_{a,hgu}$. Jednoznaczność w kogranicy daje funktorialność.
Kładziemy $\eta_a=\lambda^{K(a)}_{a,\id_{K(a)}}$. Dla $f:a\to a'$ równanie~\eqref{cat:eq:lan-action} daje
\[
 L(K(f))\eta_a=\lambda^{K(a')}_{a,K(f)}
 =\lambda^{K(a')}_{a',\id}F(f)=\eta_{a'}F(f),
\]
gdzie środkowa równość pochodzi z morfizmu $f$ w kategorii $(K\downarrow K(a'))$. Zatem $\eta$ jest naturalna.

Dla $H:\mathcal B\to\mathcal C$ i $\sigma:F\Rightarrow HK$ rodzina $H(u)\sigma_a:F(a)\to H(b)$ jest kostożkiem: jeśli $u'K(f)=u$, to
\[
 H(u')\sigma_{a'}F(f)=H(u')H(K(f))\sigma_a=H(u)\sigma_a.
\]
Istnieje jedyny $\alpha_b:L(b)\to H(b)$ z $\alpha_b\lambda^b_{a,u}=H(u)\sigma_a$. Dla $g:b\to c$ obie strony $H(g)\alpha_b$ i $\alpha_c L(g)$ po prekompozycji z $\lambda^b_{a,u}$ dają $H(gu)\sigma_a$, więc są równe. Mamy $\alpha_{K(a)}\eta_a=\sigma_a$ przez wybór $u=\id$. Jeśli $\alpha'$ również spełnia $(\alpha' K)\eta=\sigma$, to z~\eqref{cat:eq:lan-action} i naturalności $\alpha'$ wynika
\[
 \alpha'_b\lambda^b_{a,u}=\alpha'_b L(u)\eta_a
 =H(u)\alpha'_{K(a)}\eta_a=H(u)\sigma_a.
\]
Jednoznaczność w kogranicy daje $\alpha'=\alpha$. Dowodzi to własności uniwersalnej.
\end{proof}
\begin{twierdzenie}[Punktowe prawe rozszerzenie Kana]\label{cat:thm:pointwise-ran}
Jeżeli dla każdego $b$ istnieje granica diagramu $FP_b^+:(b\downarrow K)\to\mathcal C$, to istnieje prawe rozszerzenie Kana i
\[
 (\operatorname{Ran}_K F)(b)=\lim_{(a,v:b\to K(a))}F(a).
\]
W szczególności warunek zachodzi, gdy $\mathcal C$ jest zupełna.
\end{twierdzenie}
\begin{proof}
Wybieramy granicę $R(b)$ z projekcjami $\rho^b_{a,v}:R(b)\to F(a)$ spełniającymi $F(f)\rho^b_{a,v}=\rho^b_{a',v'}$ dla $K(f)v=v'$. Dla $g:b\to c$ rodzina $\rho^b_{a,wg}$ indeksowana obiektami $(a,w)$ w $(c\downarrow K)$ jest stożkiem, ponieważ $K(f)w=w'$ implikuje $K(f)wg=w'g$. Definiujemy jedyny $R(g):R(b)\to R(c)$ z
\begin{equation}\label{cat:eq:ran-action}
 \rho^c_{a,w}R(g)=\rho^b_{a,wg}.
\end{equation}
Identyczność spełnia to równanie dla $g=\id_b$. Dla $h:c\to d$ postkompozycja obu $R(hg),R(h)R(g)$ z każdą projekcją $\rho^d_{a,z}$ daje $\rho^b_{a,zhg}$, więc z jednoznaczności w granicy otrzymujemy zachowanie składania.
Definiujemy $\epsilon_a=\rho^{K(a)}_{a,\id}$. Dla $f:a\to a'$ mamy
\[
 \epsilon_{a'}R(K(f))=\rho^{K(a)}_{a',K(f)}
 =F(f)\rho^{K(a)}_{a,\id}=F(f)\epsilon_a.
\]
Środkowa równość to warunek stożka w $(K(a)\downarrow K)$.

Dla $\sigma:HK\Rightarrow F$ rodzina $\sigma_a H(v):H(b)\to F(a)$ jest stożkiem, gdyż $F(f)\sigma_a H(v)=\sigma_{a'}H(K(f)v)=\sigma_{a'}H(v')$. Istnieje jedyny $\beta_b:H(b)\to R(b)$ z $\rho^b_{a,v}\beta_b=\sigma_a H(v)$. Dla $g:b\to c$ postkompozycja $R(g)\beta_b$ oraz $\beta_c H(g)$ z każdą projekcją $\rho^c_{a,w}$ daje $\sigma_a H(wg)$, zatem $\beta$ jest naturalna. Mamy $\epsilon_a\beta_{K(a)}=\sigma_a$. Jeżeli $\beta'$ spełnia to samo równanie, to dla każdego $(a,v)$
\[
 \rho^b_{a,v}\beta'_b=\epsilon_a R(v)\beta'_b
 =\epsilon_a\beta'_{K(a)} H(v)=\sigma_a H(v).
\]
Jednoznaczność w granicy daje $\beta'=\beta$.
\end{proof}
\begin{uwaga}
Powyższe dowody podają warunki wystarczające istnienia rozszerzeń. Samo istnienie rozszerzenia w sensie własności uniwersalnej nie jest definicją rozszerzenia punktowego. Rozszerzeniami punktowymi nazywamy tu rozszerzenia obliczane za pomocą wskazanych granic lub kogranic i ich kanonicznych transformacji.
\end{uwaga}
\begin{przyklad}[Pusta kategoria przecinkowa]\label{cat:ex:kan-empty}
Jeżeli $(K\downarrow b)$ jest pusta, wartość punktowego lewego rozszerzenia w $b$ jest obiektem początkowym $\mathcal C$, o ile taki obiekt istnieje. Jeżeli $(b\downarrow K)$ jest pusta, wartość prawego rozszerzenia jest obiektem końcowym. Wynika to z definicji kogranicy i granicy pustego diagramu: rodzina składowych kostożka lub stożka jest pusta, a własność uniwersalna wymaga dokładnie jednego morfizmu w odpowiednim kierunku.
\end{przyklad}

\subsection{Rozszerzenia wzdłuż funktorów w pełni wiernych}
\begin{twierdzenie}\label{cat:thm:kan-fully-faithful}
Jeżeli $K:\mathcal A\to\mathcal B$ jest w pełni wierny, to jednostka $\eta:F\Rightarrow LK$ punktowego lewego rozszerzenia oraz kojednostka $\epsilon:RK\Rightarrow F$ punktowego prawego rozszerzenia (o ile odpowiednie rozszerzenie istnieje) są izomorfizmami naturalnymi.
\end{twierdzenie}
\begin{proof}
Obiekt $(a,\id_{K(a)})$ jest końcowy w $(K\downarrow K(a))$: morfizm z $(x,u)$ do niego to strzałka $f:x\to a$ z $K(f)=u$, istniejąca przez pełność i jedyna przez wierność. Przypomnijmy, że dla diagramu $D:\mathcal J\to\mathcal C$ z obiektem końcowym $t$ w $\mathcal J$ kogranicą jest $D(t)$. Jedynym strzałkom $j_x:x\to t$ odpowiada kostożek $D(j_x)$, ponieważ $j_yh=j_x$ dla $h:x\to y$. Każdy kostożek $(k_x:D(x)\to W)_x$ spełnia $k_x=k_tD(j_x)$, a $D(j_t)=\id$ wymusza jednoznaczność $k_t$. Stosując to do $FP^-_{K(a)}$, otrzymujemy izomorfizm $\eta_a=\lambda^{K(a)}_{a,\id}:F(a)\to L(K(a))$ z jednoznaczności kogranicy.

Dla prawego rozszerzenia $(a,\id_{K(a)})$ jest początkowy w $(K(a)\downarrow K)$: strzałka do $(x,v)$ odpowiada jedynemu $f:a\to x$ z $K(f)=v$. Gdy $s$ jest początkowy w $\mathcal J$, jedyne strzałki $j_x:s\to x$ dają stożek $D(j_x)$, a dowolny stożek $(k_x:W\to D(x))_x$ spełnia $k_x=D(j_x)k_s$. Zatem $D(s)$ jest granicą. Dostajemy izomorfizm $\epsilon_a=\rho^{K(a)}_{a,\id}$. Naturalność obu transformacji wynika z twierdzeń~\ref{cat:thm:pointwise-lan} i~\ref{cat:thm:pointwise-ran}; odwrotności składowych są naturalne przez twierdzenie~\ref{tw:nat-iso-components}.
\end{proof}
\begin{wniosek}\label{cat:cor:kan-fully-faithful}
Jeżeli $K$ jest w pełni wierny, to istniejące dla wszystkich $F$ punktowe funktory $\operatorname{Lan}_K$ i $\operatorname{Ran}_K$ są w pełni wierne (każdy przy odpowiednim założeniu istnienia).
\end{wniosek}
\begin{proof}
Jednostka sprzężenia $\operatorname{Lan}_K\dashv K^*$ i kojednostka $K^*\dashv\operatorname{Ran}_K$ są odpowiednio rodzinami $\eta_F$ i $\epsilon_F$, gdyż odpowiadają identycznościom w bijekcjach z definicji~\ref{cat:def:kan-extension}. Są izomorfizmami przez twierdzenie~\ref{cat:thm:kan-fully-faithful}. Stosujemy kryterium pełnej wierności funktora sprzężonego ze stwierdzenia~\ref{cat:prop:adj-full-faithful}.
\end{proof}
\begin{przyklad}[Porządki]\label{cat:ex:kan-poset}
Dla podzbioru $S$ zbioru uporządkowanego $P$, włożenia $K:S\hookrightarrow P$ z porządkiem indukowanym i monotonicznego $F:S\to Q$ do kraty zupełnej mamy
\[
 L(p)=\bigvee_{s\in S,\ s\le p}F(s),\qquad
 R(p)=\bigwedge_{s\in S,\ p\le s}F(s).
\]
Kategorie przecinkowe odpowiadają wskazanym podzbiorom; kostożki w $Q$ to ograniczenia górne, a stożki to ograniczenia dolne. Stąd wzory punktowe. Dla $s\in S$ obie rodziny zawierają $F(s)$, odpowiednio jako największy i najmniejszy element, więc $L(s)=F(s)=R(s)$. Jeśli monotoniczny $H:P\to Q$ przedłuża $F$, to $F(s)\le H(p)$ dla $s\le p$ oraz $H(p)\le F(s)$ dla $p\le s$. Z własności kresów $L(p)\le H(p)\le R(p)$. Dla $P=\{0<1<2\}$, $S=\{0,2\}$, $F(0)=0,F(2)=2$ i $Q=P$ otrzymujemy $L(1)=0$, $R(1)=2$. Oba rozszerzenia mogą więc być różne.
\end{przyklad}
\begin{przyklad}[Brak pełnej wierności]\label{cat:ex:kan-not-fully-faithful}
Niech $\mathcal A$ będzie dyskretną kategorią dwóch obiektów, $K:\mathcal A\to\mathbf1$ jedynym funktorem, a $F$ przypisuje obu obiektom dwuelementowe zbiory. Lewe rozszerzenie jest ich sumą rozłączną, prawe iloczynem. Jednostka składa się z włożeń, kojednostka z projekcji; żadne nie są bijekcjami. Funktor $K$ jest wierny, ale nie pełny, bo nie podnosi identyczności w $\mathbf1$ do strzałki między różnymi obiektami.
\end{przyklad}

\subsection{Rozszerzenia wzdłuż funktorów sprzężonych}
\begin{twierdzenie}\label{cat:thm:kan-along-adjoint}
Jeżeli $K\dashv T$ z jednostką $u:\id_{\mathcal A}\Rightarrow TK$ i kojednostką $v:KT\Rightarrow\id_{\mathcal B}$, to $(FT,Fu)$ jest punktowym lewym rozszerzeniem Kana $F$ wzdłuż $K$. Jeżeli $S\dashv K$ z jednostką $u':\id_{\mathcal B}\Rightarrow KS$ i kojednostką $v':SK\Rightarrow\id_{\mathcal A}$, to $(FS,Fv')$ jest punktowym prawym rozszerzeniem. Nie potrzeba zupełności ani kozupełności kategorii wartości.
\end{twierdzenie}
\begin{proof}
Dla $K\dashv T$ obiekt $(T(b),v_b)$ jest końcowy w $(K\downarrow b)$, gdyż dla $f:K(a)\to b$ sprzężenie daje dokładnie jedną $h:a\to T(b)$ z $v_bK(h)=f$. Kogranica istnieje i jest $F(T(b))$, z kostożkiem $F(h)$, przez argument z dowodu twierdzenia~\ref{cat:thm:kan-fully-faithful}. Dla $g:b\to c$ transpozycją $gf$ jest $T(g)h$, gdyż $v_cK(T(g))K(h)=gv_bK(h)=gf$. Wzór~\eqref{cat:eq:lan-action} daje zatem $L(g)=F(T(g))$. Transpozycją $\id_{K(a)}$ jest $u_a$ przez równanie trójkątne $v_{K(a)}K(u_a)=\id_{K(a)}$, więc jednostką rozszerzenia jest $Fu$.

Dla $S\dashv K$ obiekt $(S(b),u'_b)$ jest początkowy w $(b\downarrow K)$: dla $f:b\to K(a)$ istnieje jedyna $h:S(b)\to a$ z $K(h)u'_b=f$. Granica jest $F(S(b))$, z projekcjami $F(h)$. Dla $g:b\to c$ i $f:c\to K(a)$ transpozycją $fg$ jest $hS(g)$, bo $K(h)K(S(g))u'_b=K(h)u'_cg=fg$. Wzór~\eqref{cat:eq:ran-action} daje $R(g)=F(S(g))$. Transpozycją $\id_{K(a)}$ jest $v'_a$ przez równanie $K(v'_a)u'_{K(a)}=\id_{K(a)}$. Stąd kojednostka $Fv'$.
\end{proof}
\begin{przyklad}
Dla równoważności kategorii $K$ z funktorem odwrotnym $T$ z dokładnością do izomorfizmu naturalnego oba rozszerzenia są kanonicznie izomorficzne z $FT$. Istotnie równoważność można wyposażyć w oba sprzężenia (rozdział~\ref{sec:adjunctions}), a twierdzenie~\ref{cat:thm:kan-along-adjoint} stosuje się do każdego z nich. Dla włożenia $K:\mathbf1\to\mathcal B$ wskazującego obiekt początkowy $\mathcal B$ mamy $K\dashv T$ z jedynym $T:\mathcal B\to\mathbf1$, więc lewe rozszerzenie dowolnego obiektu jest funktorem stałym; dla obiektu końcowego analogicznie prawe rozszerzenie jest stałe.
\end{przyklad}

\subsection{Zachowywanie, absolutność i składanie}
\begin{definicja}\label{cat:def:kan-preserved}
Funktor $J:\mathcal C\to\mathcal D$ zachowuje lewe rozszerzenie $(L,\eta)$, jeśli $(JL,J\eta)$ jest lewym rozszerzeniem $JF$ wzdłuż $K$. Dla prawego rozszerzenia analogicznie wymagamy, aby $(JR,J\epsilon)$ było prawym rozszerzeniem. Rozszerzenie nazywamy absolutnym, gdy zachowuje je każdy funktor wychodzący z $\mathcal C$ (w rozpatrywanym uniwersum).
\end{definicja}
\begin{stwierdzenie}\label{cat:prop:kan-preservation}
Jeśli $J$ zachowuje kogranice wszystkich diagramów $FP_b^-$, to zachowuje punktowe lewe rozszerzenie Kana. Jeśli zachowuje granice wszystkich $FP_b^+$, to zachowuje punktowe prawe rozszerzenie. W szczególności lewe funktory sprzężone zachowują punktowe lewe rozszerzenia, a prawe funktory sprzężone punktowe prawe rozszerzenia.
\end{stwierdzenie}
\begin{proof}
W przypadku lewym $JL(b)$ z kostożkiem $J\lambda^b_{a,u}$ jest kogranicą $JFP_b^-$. Równanie~\eqref{cat:eq:lan-action} po zastosowaniu $J$ pokazuje, że $JL(g)$ jest dokładnie morfizmem wyznaczonym przez te kostożki; jednostką konstrukcji jest $J\eta$. Twierdzenie~\ref{cat:thm:pointwise-lan} daje tezę. W przypadku prawym $JR(b)$ z projekcjami $J\rho^b_{a,v}$ jest granicą $JFP_b^+$, a równanie~\eqref{cat:eq:ran-action} wyznacza $JR(g)$ i kojednostkę $J\epsilon$. Stosujemy twierdzenie~\ref{cat:thm:pointwise-ran}. Ostatnie zdanie wynika z twierdzenia~\ref{tw:adjoints-preserve}.
\end{proof}
\begin{wniosek}\label{cat:cor:kan-absolute}
Rozszerzenia skonstruowane w twierdzeniu~\ref{cat:thm:kan-along-adjoint} są absolutne.
\end{wniosek}
\begin{proof}
Dla dowolnego $J:\mathcal C\to\mathcal D$ to samo sprzężenie $K\dashv T$ daje lewe rozszerzenie $(JF)T$ z jednostką $(JF)u=J(Fu)$. Jest to dokładnie obraz $(FT,Fu)$. Dla $S\dashv K$ analogicznie $(JF)S$ z kojednostką $(JF)v'=J(Fv')$ jest prawym rozszerzeniem.
\end{proof}
\begin{przyklad}[Rozszerzenie nieabsolutne]\label{cat:ex:kan-not-absolute}
Dla $K$ z dyskretnej kategorii dwóch obiektów do $\mathbf1$ i $F$ o wartościach singletonowych lewe rozszerzenie ma dwa elementy. Funktor $J:\Set\to\Set$, $J(X)=X\times X$, daje zbiór czteroelementowy, natomiast lewe rozszerzenie $JF$ ma dwa elementy. Nie mogą być izomorficzne. Dla prawego rozszerzenia tego samego $F$ funktor $J'(X)=X\sqcup\{*\}$ daje dwa elementy na obrazie iloczynu, a cztery elementy w iloczynie obrazów. Zatem również prawe rozszerzenia nie zawsze są absolutne.
\end{przyklad}
\begin{twierdzenie}[Składanie rozszerzeń]\label{cat:thm:kan-composition}
Niech $K:\mathcal A\to\mathcal B$ i $M:\mathcal B\to\mathcal E$, gdzie również $\mathcal E$ jest mała. Jeśli istnieją kolejne lewe rozszerzenia $(L,\eta)$ wzdłuż $K$ i $(P,\theta)$ funktora $L$ wzdłuż $M$, to $(P,(\theta K)\eta)$ jest lewym rozszerzeniem $F$ wzdłuż $MK$. Dla kolejnych prawych rozszerzeń $(R,\epsilon)$ i $(Q,\zeta)$ otrzymujemy $(Q,\epsilon(\zeta K))$ jako prawe rozszerzenie. W szczególności, gdy istnieją, mamy kanoniczne izomorfizmy
\[
 \operatorname{Lan}_{MK}F\cong\operatorname{Lan}_M(\operatorname{Lan}_K F),
 \qquad
 \operatorname{Ran}_{MK}F\cong\operatorname{Ran}_M(\operatorname{Ran}_K F).
\]
\end{twierdzenie}
\begin{proof}
Dla $H:\mathcal E\to\mathcal C$ mamy złożenie bijekcji
\[
 \Nat(P,H)\longrightarrow\Nat(L,HM)
 \longrightarrow\Nat(F,HMK).
\]
Transformacja $\alpha$ przechodzi kolejno na $(\alpha M)\theta$ i $((\alpha M)\theta K)\eta=(\alpha MK)(\theta K)\eta$. Jest to wymagana bijekcja dla rozszerzenia wzdłuż $MK$. Dla prawego rozszerzenia składamy bijekcje $\Nat(H,Q)\to\Nat(HM,R)\to\Nat(HMK,F)$, a obraz $\beta$ to $\epsilon(\zeta K)(\beta MK)$. Jednoznaczność wynika z~\ref{cat:prop:kan-unique}. Dla trzech kolejnych funktorów oba sposoby grupowania dają tę samą złożoną transformację dzięki łączności składania; jednoznaczność porównań daje zgodność kanonicznych izomorfizmów. W przypadku funktorów rozszerzeń naturalność porównań względem $F$ wynika z równania definiującego działanie na transformacjach w dowodzie twierdzenia~\ref{cat:thm:kan-adjunction}: oba złożenia w kwadracie porównawczym mają tę samą transformację po prekompozycji z $MK$ i jednostką (odpowiednio po złożeniu z kojednostką).
\end{proof}
\begin{wniosek}\label{cat:cor:kan-limit}
Dla jedynego funktora $!:\mathcal A\to\mathbf1$ lewe rozszerzenie jest kogranicą $F$, a prawe granicą $F$; równoważność dotyczy także istnienia tych obiektów.
\end{wniosek}
\begin{proof}
Funktor $\mathbf1\to\mathcal C$ wybiera obiekt $X$, a transformacja $F\Rightarrow X!$ jest kostożkiem. Własność uniwersalna lewego rozszerzenia mówi dokładnie, że każdy kostożek rozkłada się jednoznacznie przez morfizm z jego wierzchołka. Dla transformacji $X!\Rightarrow F$ i prawego rozszerzenia otrzymujemy własność uniwersalną granicy.
\end{proof}

\subsection{Jawne modele w kategorii zbiorów}
\begin{twierdzenie}\label{cat:thm:kan-set}
Dla $F:\mathcal A\to\Set$ lewy wzór punktowy ma postać
\[
 L(b)=\left(\coprod_{a\in\Ob(\mathcal A)}
       \Hom_{\mathcal B}(K(a),b)\times F(a)\right)/{\sim},
\]
gdzie $\sim$ jest najmniejszą relacją równoważności zawierającą
\[
 (a,u'K(f),x)\sim(a',u',F(f)(x))
 \quad(f:a\to a',\ u':K(a')\to b,\ x\in F(a)).
\]
Prawy wzór punktowy ma postać
\[
 R(b)=\left\{(x_{a,v})\in\prod_{(a,v)\in(b\downarrow K)}F(a):
 F(f)(x_{a,v})=x_{a',K(f)v}\ \text{dla wszystkich }f:a\to a'\right\}.
\]
Dla $g:b\to c$ mamy $L(g)[a,u,x]=[a,gu,x]$ oraz
$(R(g)x)_{a,w}=x_{a,wg}$.
\end{twierdzenie}
\begin{proof}
Iloraz jest konstrukcją kogranicy zbiorów z twierdzenia~\ref{tw:colimits-set}, zastosowaną do $(K\downarrow b)$: obiekty indeksujące to pary $(a,u)$, a każda strzałka $f$ narzuca dokładnie wskazaną relację. Włożenia $x\mapsto[a,u,x]$ tworzą kostożek. Dla dowolnego kostożka $k_{a,u}$ funkcja $[a,u,x]\mapsto k_{a,u}(x)$ jest dobrze określona, gdyż zachowuje relacje generujące, a jednoznaczność wynika z tego, że każda klasa ma reprezentanta. Działanie $L(g)$ zachowuje relacje, bo $g(u'K(f))=(gu')K(f)$, i spełnia~\eqref{cat:eq:lan-action}. Analogicznie zbiór zgodnych rodzin jest konstrukcją granicy z twierdzenia~\ref{tw:limits-set}. Dla dowolnego stożka $k_{a,v}:W\to F(a)$ funkcja $w\mapsto(k_{a,v}(w))$ ma wartości w zbiorze zgodnych rodzin, a projekcje wyznaczają ją jednoznacznie. Działanie $R(g)$ zachowuje zgodność, bo $K(f)(wg)=(K(f)w)g$, i spełnia~\eqref{cat:eq:ran-action}. Dla pustego indeksu iloraz jest pusty, a iloczyn singletonem, zgodnie z przykładem~\ref{cat:ex:kan-empty}.
\end{proof}
\begin{stwierdzenie}\label{cat:prop:kan-representables}
Dla $a_0\in\Ob(\mathcal A)$ istnieje naturalny izomorfizm
\[
 \operatorname{Lan}_K\Hom_{\mathcal A}(a_0,-)
 \cong\Hom_{\mathcal B}(K(a_0),-).
\]
Dla presnopów (funktorów kontrawariantnych do $\Set$) analogicznie
\[
 \operatorname{Lan}_{K^{\op}}\Hom_{\mathcal A}(-,a_0)
 \cong\Hom_{\mathcal B}(-,K(a_0)).
\]
Nie trzeba zakładać pełnej wierności $K$.
\end{stwierdzenie}
\begin{proof}
W lewym wzorze z twierdzenia~\ref{cat:thm:kan-set} przypisujemy klasie $[a,u,h:a_0\to a]$ strzałkę $uK(h)$. Funkcja jest dobrze określona, ponieważ $u'K(f)K(h)=u'K(fh)$. Odwrotność wysyła $w:K(a_0)\to b$ na $[a_0,w,\id_{a_0}]$. Jedno złożenie jest identycznością, a drugie również, gdyż relacja generowana przez $h$ daje $[a_0,uK(h),\id]=[a,u,h]$. Naturalność względem $g:b\to c$ wynika z $g(uK(h))=(gu)K(h)$. Stosując wynik do $K^{\op}$, otrzymujemy wersję kontrawariantną. Jest to także konsekwencja lematu Yonedy: dla $H:\mathcal B\to\Set$ mamy $\Nat(\Hom_{\mathcal A}(a_0,-),HK)\cong H(K(a_0))\cong\Nat(\Hom_{\mathcal B}(K(a_0),-),H)$ przez twierdzenie~\ref{tw:yoneda}. Obie bijekcje są naturalne w $H$.
\end{proof}

\subsection{Przykłady algebraiczne: indukcja i koindukcja}
\begin{twierdzenie}\label{cat:thm:kan-group-modules}
Niech $\varphi:G\to H$ będzie homomorfizmem grup, $k$ ciałem, $BG,BH$ kategoriami jednoobiektowymi o morfizmach odpowiednio $G,H$ i składaniu danym przez mnożenie. Funktory $BG\to\Vect_k$ odpowiadają lewym $kG$-modułom; prekompozycja z $B\varphi$ odpowiada ograniczeniu skalarów wzdłuż $kG\to kH$. Lewe i prawe rozszerzenie Kana mają postać
\[
 \operatorname{Lan}_{B\varphi}V\cong kH\otimes_{kG}V,\qquad
 \operatorname{Ran}_{B\varphi}V\cong\Hom_{kG}(kH,V).
\]
W pierwszym wzorze $kH$ jest prawym $kG$-modułem przez $\varphi$ i lewym $kH$-modułem przez mnożenie. W drugim wzorze $kH$ jest lewym $kG$-modułem przez $\varphi$, a lewe działanie $H$ na Hom określamy przez $(h\cdot f)(x)=f(xh)$.
\end{twierdzenie}
\begin{proof}
Funktor z $BG$ wybiera przestrzeń wektorową $V$ i homomorfizm $G\to\Aut(V)$; przedłużenie liniowe daje działanie $kG$. Transformacje naturalne to dokładnie liniowe morfizmy komutujące z działaniami. To identyfikuje kategorię funktorów z kategorią modułów.

Dla $kH$-modułu $W$ i $kG$-liniowego $a:V\to W$ (z ograniczonym działaniem) definiujemy $\bar a(x\otimes v)=xa(v)$. Jest to dobrze określone, gdyż $x\varphi(g)a(v)=xa(gv)$ dla $g\in G$; przez liniowość to samo zachodzi dla $kG$. Morfizm jest $kH$-liniowy. Odwrotnie, $b:kH\otimes_{kG}V\to W$ wyznacza $a(v)=b(1\otimes v)$, które jest $kG$-liniowe, bo $1\otimes gv=\varphi(g)\otimes v$. Obie konstrukcje są odwrotne, a naturalność wynika z ich wzorów. Jednostką jest $v\mapsto1\otimes v$. Zatem definicja~\ref{cat:def:kan-extension} daje lewe rozszerzenie.

W przypadku prawym $f:kH\to V$ jest $kG$-liniowe; $h\cdot f$ również, bo $(\varphi(g)x)h=\varphi(g)(xh)$. Ponadto $(h_1\cdot(h_2\cdot f))(x)=f(xh_1h_2)$, więc wzór definiuje lewe działanie. Dla $kG$-liniowego $a:W\to V$ definiujemy $\bar a(w)(x)=a(xw)$. Jest to $kG$-liniowe w $x$, a $\bar a(hw)(x)=a(xhw)=(h\cdot\bar a(w))(x)$, więc $\bar a$ jest $kH$-liniowe. Odwrotnie, $b:W\to\Hom_{kG}(kH,V)$ daje $a(w)=b(w)(1)$; mamy $a(\varphi(g)w)=b(w)(\varphi(g))=ga(w)$. Wzory są odwrotne, ponieważ $b(xw)(1)=b(w)(x)$, i naturalne względem morfizmów $V,W$. Kojednostką jest ewaluacja w $1$. Stąd prawe rozszerzenie. Ponieważ $\Vect_k$ jest zupełna i kozupełna (iloczyny i sumy przestrzeni, jądra i kojądra różnic morfizmów, twierdzenie~\ref{cat:thm:limits-products}), oba rozszerzenia są punktowe przez jednoznaczność.
\end{proof}
\begin{przyklad}
Dla homomorfizmu $G\to\{1\}$ powyższe wzory dają odpowiednio koinwarianty $V_G=V/\operatorname{span}_k\{gv-v:g\in G,v\in V\}$ oraz inwarianty $V^G=\{v:gv=v\ \forall g\in G\}$. Istotnie w iloczynie tensorowym $1\otimes gv=1\otimes v$, a funkcja $1\otimes v\mapsto[v]$ i $[v]\mapsto1\otimes v$ są odwrotne. Morfizm $kG$-liniowy $k\to V$ jest wyznaczony przez wartość w $1$, która musi być niezmiennicza; każdy niezmienniczy wektor definiuje taki morfizm. Są to odpowiednio kogranica i granica reprezentacji jako diagramu.
\end{przyklad}
\begin{uwaga}
Wzory z twierdzenia~\ref{cat:thm:kan-group-modules} działają również dla dowolnego homomorfizmu pierścieni z jedynką $R\to S$, jako sprzężenia z funktorem ograniczenia skalarów między kategoriami modułów. Wtedy $S\otimes_R-$ i $\Hom_R(S,-)$ mają te same dowody, z $s\in S$ zamiast $h\in H$ i $r\in R$ zamiast $g\in G$. Aby traktować je jako rozszerzenia Kana w języku funktorów z jednoobiektowych kategorii pierścieni, należy wymagać addytywności: sam zwykły funktor z multiplikatywnego monoidu nie musi respektować dodawania pierścienia. Wersję rozszerzeń w kategoriach wzbogaconych omówimy osobno.
\end{uwaga}
